\documentclass[11pt]{article}

\usepackage[margin=1in,headheight=14pt]{geometry}
\usepackage[T1]{fontenc}
\usepackage{lmodern}
\usepackage{amsmath,amssymb}
\usepackage{booktabs,longtable,array}
\usepackage{enumitem}
\usepackage{fancyhdr}
\usepackage[hidelinks]{hyperref}
\usepackage{xurl}
\usepackage{microtype}

\setlength{\parindent}{0pt}
\setlength{\parskip}{0.55em}
\setlist[itemize]{leftmargin=1.5em,itemsep=0.25em,topsep=0.25em}
\setcounter{tocdepth}{2}

\pagestyle{fancy}
\fancyhf{}
\lhead{Aggression Monitor}
\rhead{Charts + V3 banner}
\cfoot{\thepage}

\newcommand{\TotalAggress}{\texttt{TotalAggress}}
\newcommand{\BuySell}{\texttt{BuySell Ratio}}
\newcommand{\Warming}{\texttt{WARMING}}
\newcommand{\Unavailable}{\texttt{UNAVAILABLE}}

\title{\textbf{Aggression Monitor}\\[0.4em]
\large What the page measures, how each chart is drawn, and how the live V3 banner is classified}
\author{}
\date{Chart implementation as of 20 August 2026; V3 banner table revised 28 September 2026; overnight window 1 October 2026; contract-roll handling specified 21 October 2026 (not yet live)}

\begin{document}
\maketitle

\begin{abstract}
This note describes the trade-to-trade aggression page and the live V3 banner (\texttt{scripts/aggression\_dashboard\_v3/}, port 5053). The page measures curve aggressive flow against a chosen Treasury-futures price. Every plot and every calculation uses a trade-to-trade clock: price is the last traded price of the selected contract, and each increment of aggression and price is the change from one print to the next. There is no one-minute resample, no mid or BBO price, and no OHLCV last-in-bucket price. The two lines labelled fair value are contemporaneous diagnostics under the same $\beta_{\mathrm{ref}}$ as the banner; they are not forecasts and they do not estimate a second beta. The live 5050 Panel~4 / Panel~5 OLS implementations are not this specification. The V3 banner decomposes the latest window of that trade-to-trade path using $\beta_{\mathrm{ref}}$ expressed in yield bp per \$1 million signed DV01. That window is 15 minutes from 18:00 ET inclusive through 03:00 ET exclusive (6pm through 2am), and 5 minutes the rest of the time. That reference is today's day-tag beta (flow-history day OLS of day bps on net DV01 for that tag only), as of 9 October 2026. The earlier last-3-hour time-of-day mean was superseded on 24 September 2026; the 10-session whole-session mean was superseded on 9 October 2026. Each ready window is one coarse STATE and one fine STATE2 from the table in this note. ``Implied passive'' is the residual price contribution left after subtracting that reference-beta aggression contribution. It is not observed resting-order identity. In the six sessions of each quarterly contract roll, the page tracks the outgoing and incoming contracts as separate series, and the curve view uses roll-only betas (section ``Contract rolls''); that section is specified, not live. This note is the specification; it does not by itself restart a live dashboard.
\end{abstract}

\tableofcontents
\clearpage

\section{Goal}

V1 answers a measurement question: how much aggressive buying and selling is occurring across the Treasury curve, and how does that stack up against the chosen contract's last-trade price path on the same prints? 

\begin{itemize}
    \item Panels 1--3 show the raw flow and price together.
    \item Panel 4 draws today's $\beta_{\mathrm{ref}}$ change-on-change line and marks large residuals as diagnostics. The former Panel 5 (the last 10 sessions under the same line) was removed from live V3 on 23 October 2026.
    \item The live V3 banner reports a coarse STATE and a fine STATE2 for the latest window (or \Warming\ / \Unavailable): 15 minutes from 18:00 ET inclusive through 03:00 ET exclusive (6pm through 2am), and 5 minutes otherwise. The states are THIN LIQ, QUIET, AGGRESS DOM, BALANCED, and EXHAUSTION, with the STATE2 names in the table below.
\end{itemize}

\section{Page, controls, and data}

\subsection{What the user selects}

Default live V3 view: price symbol TY, aggression symbols the whole curve FV+TU+TY+US, DV01 adjust on, lookback 2 hours, start time the current trading-day open (18:00 ET), auto-refresh on, live mode. $\beta_{\mathrm{ref}}$ is today's day-tag beta for those selected chips versus TY last-trade price (see ``Reference beta''). TY-only remains an optional chip set, not the default. Live V3 follows the page's price symbol, chips and DV01 setting (one live view per selection; the default always runs, at most two others), and its banner uses a 15-minute window from 18:00 ET inclusive through 03:00 ET exclusive (6pm through 2am) and a 5-minute window otherwise, with the coarse/fine table in the section ``Live V3 banner (5053)'' below.

\begin{center}
\begin{tabular}{@{}p{0.24\textwidth}p{0.70\textwidth}@{}}
\toprule
\textbf{Control} & \textbf{Effect} \\
\midrule
Price symbol & Which contract's price is drawn on the right-hand axes (FV, TU, TY, US; default TY). \\
Aggression chips & Which of those four roots contribute flow (default: all four). \\
DV01 Adjust & If on, each contract's size is multiplied by a fixed dollar DV01 before summing. \\
Lookback & How much of the built series the charts display: 2, 6, 12, or 24 hours. \\
Live vs historical & Live merges the AccessAPI buffer with a session backfill. Historical pulls a fixed window. \\
\bottomrule
\end{tabular}
\end{center}

The default start is the current trading-day open, 18:00 ET.

\subsection{Sources}

Live trades come from AccessAPI (\url{ws://100.79.30.21:8001/ws/live/stream}) into an 18-hour in-memory buffer. If that feed is stale for 30 seconds, a Databento live TBBO backup can start; it stops after 60 seconds of primary recovery. Historical and backfill trades are Databento parent-symbol trade schema pulls for ZF, ZT, ZN, and ZB. Price is the last traded price of the selected price symbol. Quote updates, BBO mids, and one-minute OHLCV are not a price source for plots or calculations. When a selected aggression symbol prints and the price symbol does not, last-trade price is held at the previous print of the price symbol.

Outside a roll window each root uses one contract per UTC date, the outright with the highest traded size that day. Inside a roll window the contracts come from the roll calendar instead (section ``Contract rolls'').



All series, increments, charts, fits, and the banner use a trade-to-trade clock. Price at $k$ is $P_k$, the last traded price of the selected price symbol at or before that print. Aggression at $k$ steps by that print's signed DV01. Book updates that are not trades do not create a point and do not move $P$ or $A$.


\section{How a trade becomes flow}

\subsection{Aggressor sign}

Each trade has a source side. The mapping is
\begin{equation}
\operatorname{sign}=
\begin{cases}
+1, & \text{side}\in\{B,\mathrm{BID}\},\\
-1, & \text{side}\in\{A,\mathrm{ASK}\},\\
0, & \text{otherwise}.
\end{cases}
\end{equation}

$+1$ is an aggressive buy (lifted the offer). $-1$ is an aggressive sell (hit the bid). Unknown sides contribute zero flow and are counted in metadata only. There is no quote-rule fallback.

\subsection{Trade-to-trade increments}

For print $k$ on a selected symbol,
\begin{align}
M_k^b &= \mathbf{1}\{\operatorname{sign}_k=+1\}\,\operatorname{size}_k\,\mathrm{DV01}_{c(k)}, \\
M_k^s &= \mathbf{1}\{\operatorname{sign}_k=-1\}\,\operatorname{size}_k\,\mathrm{DV01}_{c(k)}, \\
\Delta A_k &= M_k^b-M_k^s, \\
\Delta P_k &= P_k-P_{k-1}.
\end{align}

$M_k^b$ is Buy Aggress on that print. $M_k^s$ is Sell Aggress. $\Delta A_k$ is the signed increment over the interval from the previous print to this one. $\Delta P_k$ is the last-trade price change over the same interval.

\subsection{TotalAggress}

\TotalAggress\ is the running sum of the trade-to-trade increments $\Delta A_k$, starting at the first print after the trading-day open (18:00 ET):
\begin{equation}
A_k=\sum_{u=1}^{k}\Delta A_u.
\end{equation}

Lookback only shortens the time axis. It does not restart the sum, so the last two hours still show the session-to-date level.

\subsection{BuySell Ratio}

The raw imbalance is a signed percentage on that print, not buy volume divided by sell volume:
\begin{equation}
\mathrm{BuySellRaw}_k
=100\cdot\frac{\Delta A_k}{M_k^b+M_k^s}
\qquad\text{when }M_k^b+M_k^s\neq 0.
\end{equation}

The plotted series is an EMA of that raw ratio with span 10 prints and \texttt{ignore\_na=True}, \texttt{min\_periods=1}. A zero-sign print stays missing in the raw series and does not force the EMA to zero.

\section{The charts}

Times on the x-axis are Eastern event times of the prints. Lookback trims the time axis only. The \TotalAggress\ level stays session-to-date.

\subsection{Panel 1. TotalAggress vs Price/Yield}

Two lines, dual axis. Left: session \TotalAggress\ (blue). Right: last traded price of the selected contract (red). Each point is a print.

The dual-axis ranges are scaled so the two series occupy the same vertical box, with a 5\% margin. That visual $\alpha,\beta$ is not the fair-value OLS. It is only a Plotly range mapping.

Read: when the blue cumulative flow and the red last-trade price travel together, aggressive flow and price are moving in the same direction in the displayed window. When they diverge, last-trade price is moving for reasons other than this cumulative aggression.

\subsection{Panel 2. BuySell Ratio vs Price/Yield}

Left: the EMA \BuySell\ (green). Right: the same last-trade price (orange). Both are on the trade-to-trade index.

Read: a positive ratio means that print's signed flow was a buy; a negative ratio is a sell. Because of the EMA, a single burst is smoothed across nearby prints. 

For readability, the panel should carry this fixed caption:
\begin{quote}
Signed imbalance of that print's aggressive flow, EMA(10 prints). \(0\) = two-way nearby. \(+50\) \(\approx\) 3:1 buy-heavy. \(-50\) \(\approx\) 3:1 sell-heavy. \(\pm 100\) = only one side
\end{quote}

\begin{center}
\begin{tabular}{@{}p{0.22\textwidth}p{0.70\textwidth}@{}}
\toprule
\textbf{Reading} & \textbf{Tape} \\
\midrule
near \(0\) & buys and sells about equal \\
\(+33\) / \(-33\) & about 2:1 one side \\
\(+50\) / \(-50\) & about 3:1 \\
\(+80\) to \(+100\) (or \(-80\) to \(-100\)) & almost all lifts, or hits if negative \\
\bottomrule
\end{tabular}
\end{center}

The caption and table stay on Panel 2. They do not become banner categories.

\subsection{Panel 3. Buy Aggress, Sell Aggress, TotalAggress, and price}

Four traces on three axes:
\begin{itemize}
    \item green bars: Buy Aggress $M_k^b$;
    \item red bars: Sell Aggress, plotted as $-|M_k^s|$ so sells hang below zero;
    \item black line: \TotalAggress\ on a third axis;
    \item brown line: price on the second axis.
\end{itemize}

Read: the bars are the trade-to-trade activity; the black line is the running sum; the brown line is last-trade price. This is the only panel that shows buys and sells separately.

\subsection{Panel 4. Intraday path under $\beta_{\mathrm{ref}}$}

Same axis: actual last-trade price (solid blue) and implied price (dashed red). Grey markers sit on actual price where $|z|\geq 2$.

It uses the single $\beta_{\mathrm{ref}}$ defined below: today's day-tag beta from the flow-history day OLS. Do not regress absolute $P_k$ on absolute $A_k$.

The implied print-to-print move is $\beta_{\mathrm{ref}}\,\Delta A_k$ (in price units). The dashed line is that path accumulated from the session's first last-trade price:
\begin{equation}
\widehat{P}_k=P_1+\beta_{\mathrm{ref}}(A_k-A_1).
\end{equation}
There is no per-print intercept $\alpha(k-1)$. The residual is $e_k^{\mathrm{path}}=P_k-\widehat{P}_k$. The displayed z-score is
\begin{equation}
z_k^{\mathrm{V1,live}}=e_k^{\mathrm{path}}/s_e,
\end{equation}
where $s_e$ is the sample standard deviation of those path residuals in the currently aligned trade frame. The residual mean is not subtracted. A point is marked large when $|z_k^{\mathrm{V1,live}}|\geq 2$.

Read: the dashed line is where last-trade price would sit if each print-to-print price change equalled $\beta_{\mathrm{ref}}$ times that print's signed aggression. When solid price is above the dash, the residual is positive: price is rich relative to the reference beta. That is a description, not a prediction that price will revert.

\subsection{Panel 5 (removed)}

Live V3 no longer draws a last-10-sessions path. It was about 40\% of the first-load snapshot and was removed on 23 October 2026. Panels 1--4 and the banner are unchanged. Past states stay available by hovering any chart timestamp; the hover box looks up the banner row at or before that print.

\section{Live V3 banner (5053)}

Live V3 (\texttt{scripts/aggression\_dashboard\_v3/}, port 5053) answers one question: what kind of net flow is associated with the chosen contract's current net last-trade price action over the banner window? The path inside that window is the trade-to-trade series. Each ready window has one coarse STATE and one fine STATE2. Before classification is ready, the headline is \Warming\ or \Unavailable. A print gap longer than the window in force restarts it.

Let $k(t)$ be the last print at or before wall time $t$. The lookback $W$ is 15 minutes from 18:00 ET inclusive through 03:00 ET exclusive (6pm through 2am), and 5 minutes the rest of the time. $d_t$, $a_t$, and $b_t$ are totals over that span, directly in ticks:
\begin{align}
d_t &= \frac{P_{k(t)}-P_{k(t-W)}}{\tau_c},\\
a_t &= \frac{\beta_{\mathrm{ref}}\bigl(A_{k(t)}-A_{k(t-W)}\bigr)}{\tau_c},\\
b_t &= d_t-a_t.
\end{align}
$P$ is last traded price. $A_{k(t)}-A_{k(t-W)}$ is the sum of trade-to-trade $\Delta A$ between those two prints. $d_t$ is the net last-trade price move, $a_t$ is the modeled aggression contribution, and $b_t$ is the implied-passive / residual contribution. By construction,
\begin{equation}
\boxed{d_t=a_t+b_t.}
\end{equation}

Positive aggressive buying has a positive modeled futures-price contribution; aggressive selling has a negative one. The residual $b_t$ is whatever price movement the reference aggression model does not explain. It can reflect passive-liquidity effects, news, cash or swap markets, omitted contracts, changing impact, timing mismatch, or measurement error. It is not observed passive order flow.

\subsection{Reference beta}

$\beta_{\mathrm{ref}}$ is the only beta on this page. The banner and panel 4 both use it. It is an attribution yardstick in price points per unit of signed aggression $A$ (contracts $\times$ DV01). It is not the live 5050 Panel~4 first-lock OLS and not the live Panel~5 pooled OLS.

\paragraph{Live rule (V3, tagged day beta, 9 Oct 2026).}
Each closed session day gets exactly one tag, in this order:
\begin{enumerate}
    \item an auction day takes its tenor, \texttt{2Y}, \texttt{3Y}, \texttt{5Y}, \texttt{7Y}, \texttt{10Y} or \texttt{30Y}; a 20Y auction or two tenors on one day is \texttt{AUCTION\_OTHER};
    \item otherwise \texttt{MONTH\_END} when the day is one of the last 3 business days of the calendar month;
    \item otherwise \texttt{TIER1} when NFP, CPI, FOMC, PCE or GDP prints that day;
    \item otherwise \texttt{REGULAR}.
\end{enumerate}
An auction that is also month-end or tier-1 is an auction day, and its point trains only that auction tag. Inside a contract-roll window the curve view replaces this beta with a roll beta; the day tag is still computed, trained, and shown (section ``Contract rolls''). Tags come from the economic-events last-good file the viewers already keep (\texttt{events\_tagged\_*.json}); V3 does not query Postgres.

For each tag, the beta is the flow-history day regression, fitted only on closed days with that tag:
\[
y_d=\alpha+s\,x_d+e_d,\qquad
x_d=\sum_{\text{selected contracts}}(\text{buy}-\text{sell contracts})\times\mathrm{DV01},\qquad
y_d=\frac{P^{\text{last}}_d-P^{\text{first}}_d}{1/64}\times 0.25 .
\]
$P^{\text{first}}$ and $P^{\text{last}}$ are the first and last 1-minute trade prices of the price contract, and $y$ uses the flow-history convention of 4 ticks to a bp. The slope $s$ is reported in bp per \$1 million and converted to the attribution beta as
\[
\beta_{\mathrm{ref}}=s_{\text{bp}/\$1\text{mm}}\times 4\times\tfrac{1}{64}\big/10^6 .
\]
The page's yield display still converts $\beta_{\mathrm{ref}}$ with the price contract's DV01. A flow-history bp (4 ticks) is not a yield bp, so the two numbers differ.

V3 looks up today's tag and uses that tag's beta for the whole session. Today's move is never in today's beta. A tag with fewer than 2 days, or a nonpositive slope, shows the beta as unavailable (WARMING). It does not borrow another tag's beta.

The day points live in \texttt{workspace/analysis/aggression\_v3\_live/tag\_beta.json}, one entry per selection (price contract, aggression contracts, DV01 on/off). The first build reads V1's \texttt{flow\_history.sqlite} once with \texttt{query\_only}, in the background \texttt{v3-beta-ref} thread. On the 9 Oct 2026 build that took about 210~s for the curve view and 38~s for TY only, over 1{,}233 days from 3 Jan 2022. After each 17:00 ET close, the first page refresh after the 18:00 roll reads only the last few closed days through the \texttt{(symbol, ts\_utc)} index, adds the closed day to its tag and refits. The print path reads only the file.

\paragraph{Superseded whole-session mean (25 Sep--9 Oct 2026).}
The previous live rule took the last 10 completed Globex sessions and fitted, for each one, an OLS of last-trade price change on signed aggression change over the whole session: $\Delta P_k=\alpha_s+\beta_s\Delta A_k+e_k$, with at least 30 increments. $\beta_{\mathrm{ref}}$ was the mean of the valid $\beta_s$. The function is kept as \texttt{ensure\_session\_mean\_beta\_ref}. Archived-session pages prefer the same tagged day beta when \texttt{tag\_beta.json} is ready for that selection, and fall back to this 10-session mean only when the table is missing.

\paragraph{Superseded time-of-day mean (24 Sep 2026).}
The earlier live rule fitted only the last 3 hours before the displayed clock time on each session and recomputed every 5-minute bucket.

Month-end Treasury liquidity and balance-sheet effects, one reason for the month-end tag, are documented in \href{https://libertystreeteconomics.newyorkfed.org/2024/09/end-of-month-liquidity-in-the-treasury-market/}{Duffie and Krishnamurthy, ``End-of-Month Liquidity in the Treasury Market,'' Liberty Street Economics, 16 September 2024}.

The live 5050 dashboard is unchanged by this note: it still estimates separate session and 10-day OLS betas and does not use $\beta_{\mathrm{ref}}$ on the charts or the banner.

\subsection{Signs}

The quiet threshold is $q=1$ tick. For each of $a_t$, $d_t$, and $b_t$:
\begin{itemize}
    \item pos: the value is greater than $+q$;
    \item neg: the value is less than $-q$;
    \item small: the absolute value is at most $q$. Exactly $\pm q$ is small.
\end{itemize}

Aggression and price pick the state. The residual splits the state only when $a_t$ and $d_t$ are both pos or both neg. The identity $d_t=a_t+b_t$ leaves the 19 sign triples in the table. The states do not identify the path inside the window.

\subsection{State table}

\begin{center}
\footnotesize
\setlength{\tabcolsep}{3.5pt}
\renewcommand{\arraystretch}{1.12}
\begin{tabular}{@{}lll>{\raggedright\arraybackslash}p{0.13\textwidth}>{\raggedright\arraybackslash}p{0.27\textwidth}>{\raggedright\arraybackslash}p{0.23\textwidth}@{}}
\toprule
$a_t$ & $d_t$ & $b_t$ & STATE & STATE2 & Condition \\
\midrule
small & pos & small & THIN LIQ & SELLER STRIKE & $|a_t|\le q,\ d_t>q,\ |b_t|\le q$ \\
small & pos & pos & THIN LIQ & SELLER STRIKE & $|a_t|\le q,\ d_t>q,\ b_t>q$ \\
small & neg & small & THIN LIQ & BUYER STRIKE & $|a_t|\le q,\ d_t<-q,\ |b_t|\le q$ \\
small & neg & neg & THIN LIQ & BUYER STRIKE & $|a_t|\le q,\ d_t<-q,\ b_t<-q$ \\
small & small & small & QUIET & QUIET & $|a_t|\le q,\ |d_t|\le q,\ |b_t|\le q$ \\
small & small & pos & QUIET & QUIET & $|a_t|\le q,\ |d_t|\le q,\ b_t>q$ \\
small & small & neg & QUIET & QUIET & $|a_t|\le q,\ |d_t|\le q,\ b_t<-q$ \\
pos & pos & small & AGGRESS DOM & AGGRESS DOM & $a_t>q,\ d_t>q,\ |b_t|\le q$ \\
pos & pos & pos & AGGRESS DOM & AGGRESS DOM, THIN LIQUIDITY & $a_t>q,\ d_t>q,\ b_t>q$ \\
pos & pos & neg & BALANCED & PASSIVE SELLERS ABSORBING BUYERS & $a_t>q,\ d_t>q,\ b_t<-q$ \\
pos & neg & neg & EXHAUSTION & PASSIVE SELLERS DOMINATING BUYERS & $a_t>q,\ d_t<-q,\ b_t<-q$ \\
pos & small & small & BALANCED & PASSIVE SELLERS OFFSETTING BUYERS & $a_t>q,\ |d_t|\le q,\ |b_t|\le q$ \\
pos & small & neg & BALANCED & PASSIVE SELLERS OFFSETTING BUYERS & $a_t>q,\ |d_t|\le q,\ b_t<-q$ \\
neg & pos & pos & EXHAUSTION & PASSIVE BUYERS DOMINATING SELLERS & $a_t<-q,\ d_t>q,\ b_t>q$ \\
neg & neg & neg & AGGRESS DOM & AGGRESS DOM, THIN LIQUIDITY & $a_t<-q,\ d_t<-q,\ b_t<-q$ \\
neg & neg & pos & BALANCED & PASSIVE BUYERS ABSORBING SELLERS & $a_t<-q,\ d_t<-q,\ b_t>q$ \\
neg & neg & small & AGGRESS DOM & AGGRESS DOM & $a_t<-q,\ d_t<-q,\ |b_t|\le q$ \\
neg & small & small & BALANCED & PASSIVE BUYERS OFFSETTING SELLERS & $a_t<-q,\ |d_t|\le q,\ |b_t|\le q$ \\
neg & small & pos & BALANCED & PASSIVE BUYERS OFFSETTING SELLERS & $a_t<-q,\ |d_t|\le q,\ b_t>q$ \\
\bottomrule
\end{tabular}
\end{center}

Every ready window is one of these 19 rows and one of the 13 STATE2 names. EXHAUSTION on this banner is a flow label, separate from the move-exhaustion study. The banner is attribution only, never long or short.

\subsection{Minimal banner display}

The READY banner contains the headline, one contribution line, the reference beta, and an as-of time. For example:
\begin{quote}
\textbf{AGGRESS DOM}\\
5m price $+3.0$ ticks $\cdot$ aggression $+2.5$ $\cdot$ residual $+0.5$\\
Reference beta: REGULAR day OLS, 603 days (display units: yield bp per \$1 million DV01) $\cdot$ as of 10:44 ET
\end{quote}

At 10:44 ET the contribution line is labelled 5m. From 18:00 ET through 02:59 ET the same line is labelled 15m.

The fixed explanatory text is:
\begin{quote}
Implied passive is price movement left unexplained by measured aggression under the displayed reference beta; it is not observed passive order flow.
\end{quote}

Contribution z-scores, increasing / slowing labels, five-minute comparisons, state-duration tracking, path tags, and subwindow states are not banner categories in this version.

\subsection{Readiness and resets}

\begin{itemize}
    \item \texttt{READY}: a last trade at or before now and at or before now minus $W$, compatible units, a valid positive reference beta, and finite values satisfying $d_t=a_t+b_t$ within numerical tolerance. $W$ is 15 minutes from 18:00 ET inclusive through 03:00 ET exclusive, and 5 minutes otherwise.
    \item \Warming: no print yet at the $W$-minute-earlier endpoint since the session or configuration reset.
    \item \Unavailable: missing or stale input, no last-trade price, a unit mismatch, an invalid configuration, or a contribution-identity failure.
\end{itemize}

Reset the banner window after a session boundary, data gap, selected-price change, aggression-root change, DV01-mode change, beta-version change, or primary-contract change in a roll window. Hold coefficients, units, and configuration fixed inside a classified window.

\section{Contract rolls (specified 21 October 2026, not yet live)}

For a few sessions before each quarterly roll, volume is split between the expiring contract and the next one. This section says which contracts V3 uses in those sessions, which beta it applies, and what the page draws. Outside the roll window nothing in the sections above changes.

\subsection{Roll window}

Sessions are numbered by their 17:00 ET close date. Weekdays with no prints are not counted. The anchor is the last session of February, May, August, or November. The roll window is the anchor and the five sessions before it: six sessions.

For every root, the outgoing contract is the next quarterly expiry and the incoming contract is the one after it:
\begin{center}
\begin{tabular}{@{}lll@{}}
\toprule
\textbf{Roll month} & \textbf{Outgoing} & \textbf{Incoming} \\
\midrule
February & March (H) & June (M) \\
May & June (M) & September (U) \\
August & September (U) & December (Z) \\
November & December (Z) & March (H), next year \\
\bottomrule
\end{tabular}
\end{center}

The window comes from 65 rolls, August 2010 to August 2026, on the all-outright trade tape (\texttt{historical\_backfill/trades/}). A session counts as split when the incoming contract has 20--80\% of the outgoing-plus-incoming volume. Anchor$-5$ through the anchor is the shortest candidate window that contains at least 95\% of split sessions (95.8\%). For TY, the median incoming share is 0.1--0.2 at anchor$-5$, 0.35--0.47 at anchor$-2$, about 0.55 at anchor$-1$, and at least 0.9 from the anchor on. TY first passes 0.5 at anchor$-3$ in 2 rolls, anchor$-2$ in 16, anchor$-1$ in 17, and at the anchor in 30. In 85\% of root-rolls the other roots cross on the same session as TY.

\subsection{Two contract series}

Inside the window the highest-volume rule is not used. Each selected root has two contracts, $g\in\{\mathrm{out},\mathrm{in}\}$, and the page builds one trade-to-trade series per $g$:
\begin{align}
A_k^{g} &= \sum_{u\le k,\ u\in g}\Delta A_u, \\
P_k^{g} &= \text{last traded price of the price root's contract } g \text{ at or before print } k.
\end{align}
A print of contract $g$ on any selected root steps $A^{g}$. A print of the other contract does not touch series $g$. Each series starts at its own first print after 18:00 ET.

The crossover is the first closed session in the window on which the price root's incoming contract carries more than half of its outgoing-plus-incoming volume. The primary contract is
\begin{itemize}
    \item out, from anchor$-5$ until a crossover session has closed;
    \item in, from the session after the crossover, and on the anchor session in every roll.
\end{itemize}
The secondary contract is the other one. It is drawn only while the primary is out. Once the primary is in, the outgoing contract is dropped. After the crossover its volume is thin, and its fit on those sessions is unstable (curve slope $2.2$ bp per \$1 million, $R^2=0.08$).

\subsection{Roll beta}

Two more tags, \texttt{ROLL\_OUT} and \texttt{ROLL\_IN}, are fitted per selection, with the same form and conversion as the day-tag fit:
\[
y_d^{g}=\frac{P^{g,\text{last}}_d-P^{g,\text{first}}_d}{1/64}\times 0.25,\qquad
x_d^{g}=\sum_{\text{selected roots}}(\text{buy}-\text{sell contracts of }g)\times\mathrm{DV01}.
\]
$P^{g,\text{first}}_d$ and $P^{g,\text{last}}_d$ are the first and last prints of the price root's contract $g$ on session $d$. Each fit uses every closed roll-window session before today, from August 2010. V1's \texttt{flow\_history.sqlite} has no contract id, so these points come from the all-outright print tape (\texttt{historical\_backfill/trades/} plus the \texttt{trades\_v3\_*} cache). They are stored in \texttt{tag\_beta.json} next to the day tags. As with the day tags, fewer than 2 days or a nonpositive slope shows WARMING.

Which beta applies on a roll session:
\begin{itemize}
    \item Curve view (more than one aggression root): series $g$ uses \texttt{ROLL\_}$g$. This overrides the day tag.
    \item TY vs TY (the price root is the only aggression root): both series keep today's day-tag beta.
    \item The day tag is still computed and shown next to the roll beta.
\end{itemize}

Roll sessions stay in the day-tag training sample. Each one still trains its own day tag.

Evidence (walk-forward OLS, training strictly before each test day, tested from 2012, RMSE in flow-history bp of the day's move):
\begin{center}
\begin{tabular}{@{}llrr@{}}
\toprule
\textbf{View} & \textbf{Contract} & \textbf{Day-tag fit} & \textbf{Roll fit} \\
\midrule
Curve & out & 5.06 & 4.78 \\
Curve & in & 5.45 & 5.27 \\
Curve & front & 5.17 & 4.83 \\
TY vs TY & out & 4.71 & 4.77 \\
TY vs TY & in & 5.13 & 5.32 \\
TY vs TY & front & 4.87 & 4.87 \\
\bottomrule
\end{tabular}
\end{center}
Pooled curve slopes on roll sessions are $0.60$ (out, 95\% CI $0.46$--$0.76$) and $0.45$ (in, $0.33$--$0.59$) bp per \$1 million, against $1.01$ for REGULAR and $0.65$ for MONTH\_END. In the roll window, curve aggression moves TY about half as much as on a regular day. Splitting the roll fit before and after the crossover did not improve it.

Dropping roll sessions from the day-tag fits does not help. Walk-forward RMSE on non-roll days, curve view, with and without roll sessions: MONTH\_END $4.62\to4.91$, 5Y $4.30\to4.66$, 7Y $4.95\to4.87$, 2Y $4.85\to4.81$, REGULAR $4.92\to4.93$. 77\% of roll sessions already carry another tag: 134 month-end, 89 regular, 58 7Y, 45 5Y, 39 2Y, 16 tier-1, and 9 other auction days. Every February/May/August/November month-end falls inside the window. Source: \texttt{workspace/analysis/aggression\_v3\_roll/readout.md}.

\subsection{Page in the window}

\begin{itemize}
    \item Panels 1--3, the main lines on panel 4, and the first banner line use the primary contract.
    \item While the secondary is drawn, panel 1 adds its \TotalAggress\ and last-trade price as dashed lines. Panel 4 adds its actual and implied path under its own beta. The banner adds a second line with its own $d_t$, $a_t$, $b_t$, STATE, and STATE2. Each line is labelled with its contract code.
    \item The reference-beta line names the roll fit and the day tag.
\end{itemize}

For example, in the November 2026 roll with TYZ6 still primary:
\begin{quote}
\textbf{TYZ6 \quad AGGRESS DOM}\\
5m price $+3.0$ ticks $\cdot$ aggression $+2.5$ $\cdot$ residual $+0.5$\\
\textbf{TYH7 \quad BALANCED}\\
5m price $+1.0$ ticks $\cdot$ aggression $+2.6$ $\cdot$ residual $-1.6$\\
Reference beta: ROLL\_OUT / ROLL\_IN day OLS, 390 days (day tag MONTH\_END) $\cdot$ as of 10:44 ET
\end{quote}

\subsection{Live source of the second contract}

AccessAPI streams only the front contract (\texttt{.v.0}). The second contract's live prints need either a Databento Live subscription opened only inside roll windows, or a \texttt{.v.1} stream on AccessAPI. That choice is open. Until it is made, live mode draws the primary contract only, while archived sessions draw both from the trade cache. The causal backfill states in \texttt{historical\_backfill/} must be rebuilt under the same rule with a new label version.

\section{Parameter register}

\small
\begin{longtable}{@{}p{0.29\textwidth}p{0.24\textwidth}p{0.41\textwidth}@{}}
\toprule
\textbf{Name} & \textbf{Value} & \textbf{Role} \\
\midrule
\endfirsthead
\toprule
\textbf{Name} & \textbf{Value} & \textbf{Role} \\
\midrule
\endhead
Price default & TY last trade & Right-hand price series; not mid, not 1-minute OHLCV. \\
Aggression default & FV+TU+TY+US & Flow contributor; each print is one increment. TY-only remains an optional chip set. \\
Clock & trade-to-trade & All plots and calculations; interval is previous print to this print. \\
\TotalAggress\ & sum of print increments from 18:00 ET & Same path on charts and both fits. \\
DV01 FV/TU/TY/US & 42 / 36 / 85 / 125 & Dollars per bp; default on. \\
BuySell EMA span & 10 prints & Displayed ratio smoothing. \\
Default start & 18:00 ET trading-day open & Current Globex session. \\
Default lookback & 2 hours & Chart trim. \\
Live buffer retention & 18 hours & AccessAPI memory. \\
Backup stale / recover & 30s / 60s & Databento failover. \\
Backfill period / overlap & 240s / 900s & Session history refresh. \\
Trade-cache retention & 30 days & Historical parquet. \\
Trade price filter & $[80,150]$ & Dropped otherwise. \\
Front-month rule & max size by UTC date & One \texttt{instrument\_id} per root, outside roll windows only. \\
Roll window & anchor$-5$ .. anchor & Anchor is the last session of Feb/May/Aug/Nov; six sessions. Outgoing and incoming contracts come from the calendar, not volume. \\
Roll crossover & incoming $>50\%$ of pair & First closed window session, on the price root. Out is primary until then, and in from the next session and on the anchor. \\
Roll beta & \texttt{ROLL\_OUT}, \texttt{ROLL\_IN} day OLS & Curve view on roll sessions; TY vs TY keeps the day tag. Points from the all-outright print tape since Aug 2010. Day-tag fits keep roll sessions. \\
Roll second contract & while primary is out & Dashed lines on panels 1 and 4 plus a second banner line. Live source not yet chosen. \\
Regression form & $\Delta P_k$ on $\Delta A_k$ & Never absolute $P$ on absolute $A$. \\
$\beta_{\mathrm{ref}}$ & today's tag day-OLS beta & Only beta: banner and panel 4. File \texttt{tag\_beta.json}; print path reads the file. \\
Tag day slope & OLS of day bps on net DV01 for that tag & Flow-history cards' regression; min 2 days; nonpositive slope $\to$ WARMING. \\
Panel 4 line & $P_1+\beta_{\mathrm{ref}}(A_k-A_1)$ & Today's path under that same beta. \\
Panel 4 large residual & $|z|\geq 2$ & Grey markers. \\
Banner window & 15 min, 18:00--03:00 ET exclusive; else 5 min & Endpoints are last prints at $t$ and $t-W$. 6pm through 2am uses 15 minutes; the rest of the day uses 5. A print gap longer than $W$ restarts the window. \\
Banner impact & same $\beta_{\mathrm{ref}}$ & $a_t=\beta_{\mathrm{ref}}(A_{k(t)}-A_{k(t-W)})/\tau_c$. \\
Banner quiet threshold $q$ & 1 tick & pos above $+q$, neg below $-q$, small within $\pm q$. \\
Banner z-scores & None & Historical normalization is not used for classification. \\
Banner dynamics & None & No increasing / slowing, five-minute comparison, or state-duration category. \\
Banner completed bars & Last completed print & No one-minute forming bar; unused quotes do not create a point. \\
\bottomrule
\end{longtable}
\normalsize

\section{What we are watching}

The page is for one question: what kind of net flow is associated with the chosen contract's current net last-trade price action over the banner window (15 minutes from 18:00 ET through 03:00 ET exclusive, 5 minutes otherwise)? Every increment on that page is trade-to-trade. Read the banner states as follows.

\subsection{THIN LIQ}

STATE is THIN LIQ when aggression is small and price is pos or neg. STATE2 is SELLER STRIKE when price is pos, and BUYER STRIKE when price is neg.

\subsection{QUIET}

STATE and STATE2 are QUIET when aggression and price are both small. This does not mean the window had low gross volume, no two-way flow, or a smooth price path.

\subsection{AGGRESS DOM}

STATE is AGGRESS DOM when aggression and price are both pos, or both neg. STATE2 is AGGRESS DOM when the residual is small, and AGGRESS DOM, THIN LIQUIDITY when the residual has that same sign.

\subsection{BALANCED}

STATE is BALANCED for PASSIVE SELLERS ABSORBING BUYERS (aggression pos, price pos, residual neg), PASSIVE BUYERS ABSORBING SELLERS (aggression neg, price neg, residual pos), and the two OFFSETTING STATE2 names (aggression pos or neg while price is small).

\subsection{EXHAUSTION}

STATE is EXHAUSTION when aggression and price have opposite signs. STATE2 is PASSIVE SELLERS DOMINATING BUYERS when aggression is pos and price is neg, and PASSIVE BUYERS DOMINATING SELLERS when aggression is neg and price is pos. This label is the banner flow name, separate from the move-exhaustion study.

\subsection{Chart context outside the state labels}

Panels 1--3 can still show whether flow was concentrated, persistent, or two-way. Panel 4 can still show whether price is far from the reference line. Those chart readings are not extra banner states, and the states above do not predict the next move.

\subsection{Curve vs chosen price}

\TotalAggress\ includes FV, TU, TY, and US, but last-trade price is only the chosen contract. A wing print steps $A$ and holds $P$ at the previous TY trade. If TY is flat while the blue sum is moving, the flow may be in the wings. The working idea: a TY unexplained residual can just be other contracts dominating the sum.

These are contextual reading habits, not additional states. The banner is an attribution label, never a long or short signal.

\end{document}
